\notechapter{N-20221217}{Function Exercises: Domains, Ranges, Dirichlet Function, Parameters, and Functional Equations}

\section{Structure of the exercise sheet}

The exercises cover domains, ranges, the Dirichlet function, parameters, graph-based zero counting, and functional equations.  Each problem isolates a specific method rather than contributing to one continuous theorem.

\section{Domain problems}

For a function domain, the first rule is to list all legality conditions.  Square roots require nonnegative radicands, logarithms require positive arguments and valid bases, denominators cannot be zero, and iterated functions require every intermediate input to lie in the next domain.

For example,
\[
  y=\frac{\sqrt{x(3-x)}}{\lg (x-3)^2}
\]
requires
\[
  x(3-x)\ge0,
  \qquad
  (x-3)^2>0,
  \qquad
  \lg (x-3)^2\ne0.
\]
The first gives $0\le x\le3$.  The second excludes $x=3$.  The third excludes $(x-3)^2=1$, so $x=2$ or $4$, but only $2$ lies in the current interval.  Hence
\[
  D=[0,3)\setminus\{2\}.
\]
The important step is not the answer; it is the intersection of all restrictions.

For an expression such as
\[
  \sqrt{a^x-kb^x}
\]
in a denominator, the condition is strict:
\[
  a^x-kb^x>0.
\]
Since $b^x>0$, divide by $b^x$:
\[
  \left(\frac ab\right)^x>k.
\]
Now the inequality depends on whether $a/b>1$ or $0<a/b<1$.  This is a recurring trap in exponential and logarithmic domain problems.

\section{Floor functions}

The exercise sheet also uses the Gauss or floor function $[x]$.  The correct method is interval splitting:
\[
  [x]=m \quad\Longleftrightarrow\quad m\le x<m+1.
\]
A floor-function domain problem is therefore a finite or countable case split.  One should not treat $[x]$ as a smooth algebraic expression.

\section{Range problems}

For
\[
  y=\sqrt{x-4}+\sqrt{15-3x},
\]
the domain is $4\le x\le5$.  Put $u=x-4$, so $0\le u\le1$ and
\[
  y=\sqrt u+\sqrt{3(1-u)}.
\]
By Cauchy--Schwarz,
\[
  y^2\le (1+3)(u+1-u)=4,
\]
so $y\le2$.  At the endpoints the smaller value is $1$, and the maximum is achieved when the Cauchy equality condition holds.  Thus the range is $[1,2]$.

For
\[
  y=x+\sqrt{1-2x},
\]
let $t=\sqrt{1-2x}\ge0$.  Then
\[
  x=\frac{1-t^2}{2},\qquad
  y=-\frac12(t-1)^2+1.
\]
Hence the range is
\[
  (-\infty,1].
\]
This substitution method is the natural way to remove a radical while preserving the domain.

\section{Dirichlet function}

The Dirichlet function is
\[
  D(x)=\begin{cases}
  1,&x\in\mathbb Q,\\
  0,&x\in\mathbb R\setminus\mathbb Q.
  \end{cases}
\]
It is bounded, because $0\le D(x)\le1$.  It is not monotone on any interval, because every interval contains both rational and irrational numbers.  It has every nonzero rational period $T$, since
\[
  x\in\mathbb Q \Longleftrightarrow x+T\in\mathbb Q
\]
for $T\in\mathbb Q$.  It has no irrational period, because adding an irrational changes rationality.

This exercise is included because it teaches that functions can be extremely discontinuous while still having simple algebraic properties.

\section{Graph transformations and zero counting}

The problem involving $f$ and
\[
  g(x)=b-f(2-x)
\]
asks for the values of $b$ for which $f(x)-g(x)$ has four zeros.  The transformation $x\mapsto2-x$ reflects the graph across the vertical line $x=1$; then $b-$ reflects vertically and shifts.  The number of zeros of $f-g$ is the number of intersections of two transformed graphs.

The method is graphical but not vague.  The number of intersections changes only when a moving graph becomes tangent, passes through a vertex/corner, or crosses an endpoint.  Thus parameter intervals are found by checking these critical positions.

\section{Logarithmic parameter equation}

For an equation such as
\[
  \ln(ax)=2\ln(x+1),
\]
first impose the domain:
\[
  ax>0,
  \qquad x+1>0.
\]
Then combine logarithms:
\[
  \ln(ax)=\ln((x+1)^2)
\]
so
\[
  ax=(x+1)^2.
\]
This becomes a quadratic in $x$.  Having exactly one real solution means a discriminant or tangency condition, but the domain must still be checked.  The sheet is training this two-step habit: algebraic transformation plus domain verification.

\section{Functional equations with period-like recurrence}

Suppose $f$ is defined on $[0,1)$ by
\[
  f(x)=2x-x^2
\]
or another explicit expression, and for all real $x$ satisfies
\[
  f(x)+f(x+1)=1.
\]
Then values outside $[0,1)$ are reduced by repeatedly using
\[
  f(x+1)=1-f(x).
\]
If $a=\log_2 3$, to compute $f(a)+f(2a)+f(3a)$, first determine the integer intervals containing $a,2a,3a$, reduce their fractional parts into $[0,1)$, and then apply the recurrence each time.

This is the same idea as reducing modulo a period, except that shifting by $1$ may flip the value rather than preserve it.

\section{A reciprocal functional equation}

The sheet also has a functional equation of the form
\[
  f(x)=f(1)x+f(2)\frac1x-1
\]
for nonzero $x$.  On $(0,+\infty)$, this becomes a problem about minimizing
\[
  Ax+\frac Bx-1.
\]
The AM--GM inequality gives
\[
  Ax+\frac Bx\ge2\sqrt{AB}
\]
when $A,B>0$.  The equality condition determines the point where the minimum occurs.  This is a typical example where a functional equation reduces to a standard inequality.

\section{Quadratic iteration problem}

For a quadratic
\[
  f(x)=x^2+2ax+8,
\]
conditions involving both $f(x)\le0$ and $f(f(x))\le8$ should be treated by introducing the image variable $u=f(x)$.  Then the second condition becomes $f(u)\le8$.  The problem is about how the range of the first inequality is mapped by the quadratic again.  Completing the square and tracking intervals is safer than expanding $f(f(x))$ blindly.

\section{Max and absolute value}

The exercise about
\[
  M(x)=\max\{f(x),g(x)\}
\]
uses the identity
\[
  \max\{u,v\}=\frac{u+v+|u-v|}{2}.
\]
Thus maximum and minimum operations can be encoded using absolute value.  This is useful when deciding whether a function built by max/min remains elementary.


\section{Parameter problems and critical values}

When a parameter changes the number of solutions, the change usually happens at a critical value: a tangent contact, an endpoint contact, or a collision of two roots.  This is why discriminants and graph tangencies appear so often.

For a quadratic equation depending on a parameter, the condition for exactly one real root is usually
\[
  \Delta=0.
\]
For a graphical equation, the same condition appears geometrically as tangency.  The algebraic and geometric languages are two versions of the same idea.

\section{Functional equations as reduction rules}

A functional equation such as
\[
  f(x+1)=1-f(x)
\]
should be used as a reduction rule.  It tells us how to move an input back to a fundamental interval.  After two shifts,
\[
  f(x+2)=1-f(x+1)=f(x),
\]
so the function becomes periodic with period \(2\).  This kind of hidden periodicity is easy to miss if one treats the equation only algebraically.

\section{What the pattern suggests}

The page is not just a set of answers.  It teaches a toolbox: domains are legality intersections; ranges use substitution and inequalities; floor functions require interval splitting; Dirichlet's function uses density of rationals and irrationals; parameter problems are tangency/discriminant problems; functional equations reduce values to a fundamental interval; max and min can be represented by absolute values.  The common theme is that a function is not merely a formula but a domain, graph, transformation rule, and set of constraints.

\section{Domains as constraint intersections}

A domain problem is a constraint-satisfaction problem.  Each operation contributes a legality condition: square roots impose inequalities, logarithms impose positivity, denominators impose nonvanishing, and compositions impose intermediate-domain conditions.  The domain is the intersection of all these constraints.

This viewpoint prevents a common error: simplifying an expression before recording its original restrictions.  For example, an expression with a denominator may algebraically simplify, but the excluded zeros of the original denominator still matter unless the problem explicitly defines the simplified expression as a new function.

\section{Dirichlet function through density and measure}

The Dirichlet function
\[
  D(x)=\mathbf 1_{\mathbb Q}(x)
\]
is discontinuous at every real number.  At any point \(a\), every neighborhood contains rationals and irrationals, so along rational sequences \(D\) tends to \(1\), while along irrational sequences it tends to \(0\).

It is not Riemann integrable on any interval, because every subinterval has supremum \(1\) and infimum \(0\).  But it is Lebesgue integrable, and its integral is \(0\), since \(\mathbb Q\) has measure zero:
\[
  \int_0^1 \mathbf 1_{\mathbb Q}(x)\,dx=0.
\]
Thus the same function separates two integration theories: Riemann integration sees oscillation on every interval; Lebesgue integration sees that the exceptional set is small in measure.

\section{Thomae's function as a nearby comparison}

A useful comparison is Thomae's function
\[
  T(x)=
  \begin{cases}
  \frac1q,&x=\frac pq\text{ in lowest terms},\\
  0,&x\notin\mathbb Q.
  \end{cases}
\]
It is continuous at every irrational point and discontinuous at every rational point.  Unlike the Dirichlet function, it is Riemann integrable with integral \(0\).

This comparison shows that dense discontinuities alone do not determine integrability.  The size and height distribution of the discontinuities matter.

\section{Parameter problems as bifurcation diagrams}

When a parameter changes the number of solutions, the change occurs at critical values.  Algebraically, these are discriminant-zero cases, endpoint hits, or degeneracies.  Geometrically, they are tangencies or changes in intersection pattern.

For a family \(F(x,b)=0\), a typical critical point satisfies
\[
  F(x,b)=0,\qquad \frac{\partial F}{\partial x}(x,b)=0.
\]
The first equation says the graph intersects the axis; the second says the intersection is tangent, so two roots can merge or disappear.  This is the elementary version of a bifurcation diagram.

\section{Functional equations as group actions on the domain}

A relation such as
\[
  f(x+1)=1-f(x)
\]
should be read as an action of the translation group on the input line, together with an action on values.  Applying the shift twice gives
\[
  f(x+2)=f(x),
\]
so the function descends to a period-two structure after accounting for the value flip.

The interval \([0,1)\) is a fundamental domain: every real input can be moved into it by integer shifts.  The functional equation tells how the value changes during that reduction.  This is the same logic behind periodic functions, modular arithmetic, and covering spaces.

\section{Max, min, and nonsmooth points}

The identity
\[
  \max\{u,v\}=\frac{u+v+|u-v|}{2}
\]
turns a maximum into an expression involving absolute value.  It also explains where nonsmoothness can occur: at points where \(u=v\), the absolute value has a corner unless the crossing is tangential in a compatible way.

In optimization language, maximum functions are piecewise-smooth.  Their derivatives are governed by the active branch.  At a switching point, one may need one-sided derivatives or subgradients rather than an ordinary derivative.

\section{Ranges as images and compactness}

A range problem asks for the image \(f(D)\).  If \(D\) is compact and \(f\) is continuous, then \(f(D)\) is compact.  In \(\mathbb R\), that means closed and bounded.  This is why many range problems on closed intervals reduce to finding extrema.

If the domain is not compact, endpoints may be missing or infinity may be approached.  Then the range must be described with open or half-open endpoints.  This is not a cosmetic detail: it records whether a value is achieved or only approached.

\section{Function exercises through topology and dynamics}

These function exercises are training a structural habit.  A function is not just a formula; it is a formula together with a domain, a range, transformations, constraints, and possible exceptional sets.  Domain questions are intersections of constraints, range questions are image problems, parameter questions track critical values, and functional equations reduce inputs by a symmetry or group action.
